Higgs 机制改变了质量本征态，却没有改变重整化的逻辑：裸理论仍由规范不变的场与耦合定义，反项在选定真空附近旋转到质量本征基底后，再由实验输入条件固定。以维数正则化 $d=4-2\epsilon$ 为例，可写
\begin{align*}
\mathcal W_{0\mu}^I&=Z_W^{1/2}\mathcal W_\mu^I,
&\mathcal B_{0\mu}&=Z_B^{1/2}\mathcal B_\mu,
&\Phi_0&=Z_\Phi^{1/2}\Phi,\\
g_0&=\mu_{\rm DR}^{\epsilon}(g+\delta g),
&g_0'&=\mu_{\rm DR}^{\epsilon}(g'+\delta g'),
&v_0&=v+\delta v.
\end{align*}
费米子、Yukawa 耦合和 Higgs 自耦合也作同类分解。反项的允许结构受 Slavnov--Taylor 恒等式约束，其系数由选定的重整化条件和裸参数分解确定。

在壳方案把物理质量定义为传播子的实极点，并把极点留数归一为一。对 $V=W,Z$ 的横向自能，这两个条件给出
\begin{equation}
\delta m_V^2c^2
=\Re\Sigma_{VV,T}^{\rm loop}(m_V^2c^2),
\qquad
\delta Z_V
=-\Re\Sigma_{VV,T}^{{\rm loop}\,'}(m_V^2c^2).
\label{eq:ew-onshell-mass-field-ct-dp11}
\end{equation}
因此质量反项由自能在极点处的值固定，场强反项由极点处的斜率固定。光子则由未破缺的 $U(1)_{\rm em}$ 约束：重整化后的横向逆传播子在 $p^2=0$ 保持零本征值，电荷由 Thomson 极限定义；光子--$Z$ 混合反项必须与电荷和场强反项共同满足电弱 Ward/Slavnov--Taylor 关系。

若采用在壳弱混合角
\begin{equation}
s_W^2\equiv1-\frac{m_W^2}{m_Z^2},
\label{eq:onshell-sw-definition-dp11}
\end{equation}
则 $s_W$ 的反项由两个质量反项联动确定：
\begin{equation}
\boxed{
\frac{\delta s_W}{s_W}
=-\frac{c_W^2}{2s_W^2}
\left(\frac{\delta m_W^2}{m_W^2}
-\frac{\delta m_Z^2}{m_Z^2}\right).}
\label{eq:sw-counterterm-massrelation-dp11}
\end{equation}
这使弱混合角的量子修正直接连接到 $W/Z$ 两点函数。

树级关系 $g=g_{\rm em}/s_W$ 与\eqnrefs{eq:gf-matching-dp11} 给出
\begin{equation*}
M_W^2s_W^2=\frac{\pi\alpha}{\sqrt2G_F^{(E)}},
\qquad M_W\equiv m_Wc^2.
\end{equation*}
再使用\eqnrefs{eq:onshell-sw-definition-dp11}，得到三个常用输入 $(\alpha,G_F,M_Z)$ 对 $M_W$ 的树级预测
\begin{equation}
\boxed{
M_W^2\left(1-\frac{M_W^2}{M_Z^2}\right)
=\frac{\pi\alpha}{\sqrt2G_F^{(E)}}.}
\label{eq:mw-mz-alpha-gf-tree-dp11}
\end{equation}
这条关系说明精密电弱检验的基本思想：选取一组高精度输入后，其余可观测量成为量子修正和新物理的预测量。

一圈 $\mu$ 子衰变同时包含 $W$ 自能、两个重整化带电流顶角、盒图和参数反项。把这些按同一方案组合，可定义规范无关的 $\Delta r$：
\begin{equation}
\boxed{
\frac{G_F^{(E)}}{\sqrt2}
=\frac{\pi\alpha}{2s_W^2M_W^2}(1+\Delta r).}
\label{eq:delta-r-definition-dp11}
\end{equation}
单独的自能、顶角和盒图一般依赖规范和重整化方案；只有进入同一个物理输入关系后的组合才具有可观测意义。因而 $\Delta r$ 不是某一张图的“修正因子”，而是给定输入方案下完整的一圈匹配量。

在讨论质量分裂的一圈效应以前，先说明这里要检验的近似对称性。若暂时忽略超荷耦合以及同一弱二重态内部的 Yukawa 质量分裂，Higgs 标量势及其真空期望值具有比显式规范群更大的近似整体旋转对称性；电弱破缺以后，其中一个对角的 $SU(2)$ 仍保持真空不变。这个残余近似对称性把三个弱同位旋方向联系起来，并在树级给出 $m_W^2=m_Z^2c_W^2$，等价地使由中性流与带电流强度构造的参数满足 $\rho=1$。这一保护树级质量关系的近似整体对称性称为\emph{守护对称性}（custodial symmetry）。同一 $SU(2)_L$ 双重态内的质量分裂会破坏这种关系，因此最灵敏的一圈量不是单个 $W$ 或 $Z$ 自能，而是二者的差：
\begin{equation}
\Delta\rho
\equiv
\frac{\widehat\Sigma_{ZZ,T}(0)}{m_Z^2c^2}
-\frac{\widehat\Sigma_{WW,T}(0)}{m_W^2c^2}.
\label{eq:delta-rho-definition-dp11}
\end{equation}
若双重态的两个费米子具有颜色多重度 $N_c$ 和质量能量 $M_i=m_ic^2$，一圈结果为
\begin{equation}
\boxed{
\Delta\rho_f
=\frac{N_cG_F^{(E)}}{8\pi^2\sqrt2}
\left[
M_1^2+M_2^2
-\frac{2M_1^2M_2^2}{M_1^2-M_2^2}
\ln\frac{M_1^2}{M_2^2}
\right].}
\label{eq:delta-rho-doublet-dp11}
\end{equation}
这个有限结果来自带电流与中性流自能的守护对称组合：共同的紫外局域部分在组合成 $\Delta\rho$ 时已经相消，留下的函数只测量同一弱二重态内部的质量分裂。简并极限 $M_1=M_2$ 恢复这一守护对称性并使 $\Delta\rho_f\to0$。对 $M_t\gg M_b$，则
\begin{equation}
\boxed{
\Delta\rho_{tb}
=\frac{3G_F^{(E)}M_t^2}{8\pi^2\sqrt2}
\left[
1+\mathcal O\!\left(
\frac{M_b^2}{M_t^2}\ln\frac{M_t^2}{M_b^2}
\right)
\right].}
\label{eq:delta-rho-top-dp11}
\end{equation}
重顶夸克项随 $M_t^2$ 增长而不是只含对数，是因为顶夸克质量本身来自 $y_t v_E$；取大 $M_t$ 同时增大了与 Higgs/Goldstone 扇区的 Yukawa 耦合。$\Delta\rho$ 通过\eqnrefs{eq:sw-counterterm-massrelation-dp11} 进入 $\Delta r$，所以精密 $M_W$、$Z$ 极点和低能中性流测量都对守护对称性的破缺非常敏感。
